% Methods and Results for the IPA analysis: paste where the sections belong.
% The text refers to the flowchart as Figure~\ref{fig:ipa-flowchart}.
% The Results follows the report's wording of Nov 13th, 2026, as revised after the notebooks' review. Each
% paragraph cites its replication notebook in notebooks/ (entries in notebooks.bib: add \bibliographystyle{plain}
% and \bibliography{notebooks} before \end{document}, and \usepackage{url} or hyperref for the links).

\section{Methods}
We summarize each step of Figure~\ref{fig:ipa-flowchart} in the figure's row order, and the Definitions section defines the terms. Ticket counts are unplanned tickets unless stated otherwise. We rebuilt the midterm analysis's case table from IPA's service case export, kept its rule for unplanned tickets, and linked the tickets by exact serial number to the clearView telemetry database and to IPA's installed-base list, which supplies the fleets behind every rate per installed machine-month. IPA's answers on Oct 13th set how we read ticket times and fields. Its two questions on slides 4 and 5 set the follow-ups, and its requests on Oct 30th set the Electronics case list and a redrawn chart for slide 4a.

\smallskip\noindent\textbf{Customer fleet ownership.}
We counted the telemetry devices that ticket serials reach and compared accounts with sites.

\smallskip\noindent\textbf{High-ticket machines.}
We flagged rate outliers and chronic outliers [one-sided exact Poisson test; Benjamini--Hochberg], compared ticket concentration with chance [Poisson simulation], re-tested the rate outliers allowing for spread between machines [negative binomial], and compared the rise in non-success codes before tickets on chronic outliers and other machines.

\smallskip\noindent\textbf{Repeat tickets.}
We measured the share of eligible tickets followed by another ticket on the same machine within 30 days and within 7 days, and the 30-day share without monitoring tickets. With the follow-ups, we compared the first two with random timing [exact calculation; Monte Carlo check].

\smallskip\noindent\textbf{Rates by model and age.}
We computed rates per 100 installed machine-months by model, age, and service status and compared each model with the fleet raw and age for age [indirect standardization; Poisson main-effects model; iterative proportional fitting; cluster bootstrap].

\smallskip\noindent\textbf{Customers vs their fleets.}
We compared each customer's tickets with its fleet's expected tickets [one-sided exact Poisson test; Benjamini--Hochberg; overdispersion-adjusted funnel plot] and compared ticket concentration with chance [Poisson simulation].

\smallskip\noindent\textbf{Persistence and outliers.}
We split the period into two halves, refitted expected tickets in each, and compared the halves [Spearman rank correlation; stratified permutation test]. We flagged consistent outliers [one-sided Poisson tests; Benjamini--Hochberg] and machines above their own expected tickets beyond chance [one-sided exact Poisson test; Benjamini--Hochberg; negative binomial].

\smallskip\noindent\textbf{Ticket clock and fields.}
We checked whether opening times step backwards in case-number order on the export's clock and on hospital-local clocks. We compared how often the issue fields are filled on open tickets and on August 2026 tickets with status Trouble Closed.

\smallskip\noindent\textbf{First look at codes.}
For the paired tickets, we compared the last 10 codes before each ticket with those in its comparison window by code group, with no test.

\smallskip\noindent\textbf{Association results.}
Following a plan and a Category-to-code crosswalk fixed in advance, we compared the odds that a ticket's own codes appear in its final 7 days rather than in the final 7 days of its comparison window with the same odds for other faults' tickets. The ratio is adjusted for machine type and quarter and pooled over the code sets that could be tested [Mantel--Haenszel odds ratio; one-sided permutation test; cluster bootstrap]. We then ran planned sensitivity checks.

\smallskip\noindent\textbf{Request types by age.}
To answer whether young machines' extra tickets are different kinds of problem, we compared rates on young and older machines by request group, all models together [cluster bootstrap], and compared the mix of request groups within model [permutation test; Benjamini--Hochberg; Mantel--Haenszel share ratio]. We split the extra tickets into within-model and model-mix parts and, after seeing the data, grouped machines by install year.

\smallskip\noindent\textbf{Repeat-contact pairings.}
To answer what follows what in repeat contacts, we classified each repeat and compared same-Category shares by the gap between a machine's tickets. We compared each model and telemetry type that has at least 100 eligible tickets with random timing, relative to the fleet [cluster bootstrap; Benjamini--Hochberg].

\smallskip\noindent\textbf{Electronics cases.}
As a descriptive step outside the plan, we compared Electronics and its Categories on young and older machines, all models together, and examined Touchscreen tickets by Sub Category, model, and install date. We built a case list, which an independent rebuild checked.

\smallskip\noindent\textbf{Verification.}
We recomputed every case-export number of the midterm analysis three independent ways and re-ran its telemetry results on the full database. The association study and the answers to IPA's two questions followed plans fixed in advance, with later changes recorded as dated amendments. We checked each analysis with a review or an independent rebuild, written from the definitions or plan without reading the analysis code. The exception is the Oct 14th first look at codes, which the association study later reproduced.

\section{Results}
\smallskip\noindent\textbf{Customer fleet ownership:}
Ticket serials reach 3,041 of 7,183 devices (42\%), and accounts are individual facilities (845 accounts at 725--837 sites), so fleets come from the installed-base list~\cite{nb01}.

\smallskip\noindent\textbf{High-ticket machines:}
We found 113 chronic outliers at 82 customers. The top 5\% of the 6,621 telemetry devices hold 32.9\% of their tickets vs. 22.2\% by chance. Allowing for spread between machines, only 2 rate outliers still stand out, so high-ticket machines form a continuum. Non-success codes rise before their tickets no more than on other machines~\cite{nb02}.

\smallskip\noindent\textbf{Repeat tickets:}
Of 6,665 eligible tickets, 23.6\% repeat within 30 days and 9.5\% within 7 days, 1.53 and 2.30 times random timing's shares. Without monitoring tickets, the 30-day share is 22.0\%~\cite{nb03}.

\smallskip\noindent\textbf{Rates by model and age:}
The fleet runs at 6.2 unplanned tickets per 100 installed machine-months. The three busiest models (LVD-Remote-2.0, LVRR-WP-2.0, and MV-2.0) hold 42\% of machine-months and 51\% of tickets, partly because they are younger. scrubBank 480 runs 2.6 times the fleet age for age. Machines under 4 years run about 8 per 100 vs. 2.5 at 10 years or more, but the 10+ rate overstates the drop with age: 83 of its 164 machines with a ticket are Not Serviceable, and removed machines are missing from the list~\cite{nb04}.

\smallskip\noindent\textbf{Customers vs. their fleets:}
63 customers exceed their fleets beyond chance (1,732 tickets vs. 747 expected), and 1 lies outside the normal spread. The top 10\% of customers hold 47\% of tickets vs. 39\% from their fleets plus chance~\cite{nb04}.

\smallskip\noindent\textbf{Persistence and outliers:}
Customers above their fleets in the first half tend to stay above them in the second (rank correlation 0.46 vs. 0.04 when shuffled), and machines with 3 or more first-half tickets had 2.6 times their expected tickets in the second half. 18 customers and 36 machines are consistent outliers. Of the 106 chronic outliers on IPA's list, 93 stay flagged, and no flagged machine stands out as a separate group~\cite{nb04}.

\smallskip\noindent\textbf{Ticket clock and fields:}
Opening times never step backwards in case-number order (0 of 7,470 steps) but would 1,295--1,312 times on hospital-local clocks, so case times use one U.S. Eastern clock, as IPA confirmed. Trouble Code and Category are filled on 1 of 64 open tickets and on nearly all of the 461 August 2026 tickets of any kind with status Trouble Closed. Category is set near close, and the export holds only each field's final value~\cite{nb05}.

\smallskip\noindent\textbf{First look at codes:}
Door codes appear before 68\% of the 3,259 paired tickets vs. 61\% of their comparison windows, and electronics codes before 20\% vs. 13\%~\cite{nb06}.

\smallskip\noindent\textbf{Association results:}
The odds that a ticket's own codes appear in its final 7 days rather than in its comparison window's final 7 days are 3.2 times those of other faults' tickets (95\% interval 2.3--4.8). By code set, the ratios are rough and do not rank the four sets with ratios above 1: Main Door is borderline, and Computer Board was not shown to be specific. At 7 days, the pooled result holds in every sensitivity check run so far that gives one, but it rests almost entirely on scrubex machines and is provisional~\cite{nb06}.

\smallskip\noindent\textbf{Request types by age:}
Young machines run 8.0 tickets per 100 machine-months vs. 4.8 for older ones (1.66 times), with higher rates in nearly every request group. Within model, Electronics' share is 1.31 times as large on young machines (mostly Touchscreen) and Mechanical's 0.70 times. 80\% of the 1,447 extra tickets come from young machines running above older ones of the same model. Touchscreen tickets cluster on 2022--23 installs (1.54 per 100, 3.9 times the 2024--26 rate)~\cite{nb07}.

\smallskip\noindent\textbf{Repeat-contact pairings:}
Of the 1,574 repeats, 31\% are the same Category, 10\% the same request group with another Category, and 32\% a different group. The other 28\% involve monitoring, No Trouble Found, or a blank Category. Over all pairs of a machine's tickets, 46\% of pairs a day or less apart share a Category vs. 27\% of pairs over six months apart. No model or telemetry type with 100 or more eligible tickets clearly stands out from the fleet~\cite{nb03}.

\smallskip\noindent\textbf{Electronics cases:}
Per machine-month, Electronics runs 1.71 times as often on young machines, near the 1.66 for all tickets. Touchscreen stands out at 4.2 times: of its 411 tickets on young machines, 344 are connector reseats or replacements, 402 are on the four 2.0 models (LVD-Remote-2.0, MV-2.0, LVRR-WP-2.0, and SV-48-2.0), and 315 are on 2022--23 installs. Its rate steps from 0.36 to 1.61 per 100 between Jan--Apr and May--Aug 2022 installs, which are nearly the same age, so the step points to a production, part, or installation change. The case list has all 1,605 unplanned Electronics cases on the fleet machines behind slide 4~\cite{nb07}.

\smallskip\noindent\textbf{Verification:}
Every case-export number reproduced exactly, and on the full database 65 of 66 checks passed; the one failure was a sort-order bug in our comparison script. Independent rebuilds matched the 113 chronic outliers, the fleet rates before a statistical review's six changes, all 33,515 compared values of the association study, the follow-ups' headline numbers, and every sheet of the case list. Reviews found that the chronic-outlier count runs from 62 to 133 with the statistical cut-off. The association results' 45 review points were all upheld and applied, and no headline result changed~\cite{nb08}.
